\documentclass[11pt]{article}
\usepackage[spanish]{babel}
\usepackage{amsmath}
\usepackage{amssymb}
\usepackage[margin=0.75in]{geometry} % Márgenes más pequeños
\usepackage[T1]{fontenc}
\usepackage[utf8]{inputenc}
\usepackage{multicol} % Para usar dos columnas
\usepackage[most]{tcolorbox} % Para las cajas de colores
\usepackage{fancyhdr} % <<--- PAQUETE AÑADIDO (soluciona los errores)
\usepackage{enumitem} % <<--- PAQUETE AÑADIDO (para compactar listas)

% --- Configuración de las cajas de TColorBox ---
\tcbset{
    colback=blue!5!white, % Color de fondo de la caja
    colframe=blue!75!black, % Color del marco
    fonttitle=\bfseries,
    arc=2mm, % Esquinas redondeadas
    boxsep=5pt,
    left=6pt,
    right=6pt,
    top=6pt,
    bottom=6pt
}


\title{Formulario Completo de Física Teórica}
\author{Felipe Colli \thanks{AFTIN y Profesor Paul Cáceres}}
\date{2025}

% --- Configuración del Pie de Página ---
\pagestyle{fancy} % Activa el estilo fancy para todo el documento
\fancyhf{} % Limpia encabezado y pie de página
\fancyfoot[C]{\small \textcopyright{} 2025, Felipe Colli. Todos los derechos reservados.}
\fancyfoot[R]{\thepage}
\renewcommand{\headrulewidth}{0pt} % Sin línea en el encabezado
\renewcommand{\footrulewidth}{0.4pt} % Línea en el pie de página

\begin{document}

\maketitle
\thispagestyle{fancy} % Se asegura que la primera página también use el estilo fancy

\begin{multicols}{2} % Inicia el entorno de dos columnas
\setlist{nosep} % Reduce el espacio vertical en todas las listas de aquí en adelante

%===========================================================
% CINEMÁTICA LINEAL
%===========================================================
\begin{tcolorbox}[title=Cinemática Lineal]
\subsection*{Movimiento Rectilíneo Uniforme (MRU)}
\begin{itemize}
    \item $v = \frac{\Delta x}{\Delta t}$
    \item $x(t) = x_i + v \cdot t$
\end{itemize}

\subsection*{Movimiento Rectilíneo Uniformemente Acelerado (MRUA)}
\begin{itemize}
    \item $v_f = v_i + a \cdot t$
    \item $x(t) = x_i + v_i t + \frac{1}{2}at^2$
    \item $v_f^2 = v_i^2 + 2a\Delta x$
\end{itemize}

\subsubsection*{Definiciones}
\begin{itemize}
    \item $x$: Posición [m]
    \item $v$: Velocidad [m/s]
    \item $a$: Aceleración [m/s$^2$]
    \item $t$: Tiempo [s]
\end{itemize}
\end{tcolorbox}


%===========================================================
% MOVIMIENTO CIRCULAR
%===========================================================
\begin{tcolorbox}[title=Movimiento Circular]
\subsection*{Relaciones Lineal-Angular}
\begin{itemize}
    \item $s = \theta \cdot r$ (Arco)
    \item $v_t = \omega \cdot r$ (Velocidad tangencial)
    \item $a_t = \alpha \cdot r$ (Aceleración tangencial)
\end{itemize}

\subsection*{Movimiento Circular Uniforme (MCU)}
\begin{itemize}
    \item $a_c = \frac{v_t^2}{r} = \omega^2 r$ (Acel. centrípeta)
    \item $T = \frac{2\pi}{\omega} = \frac{1}{f}$ (Periodo)
    \item $f = \frac{1}{T}$ (Frecuencia)
\end{itemize}

\subsection*{Movimiento Circular Uniformemente Variado (MCUV)}
\begin{itemize}
    \item $\omega_f = \omega_i + \alpha t$
    \item $\theta(t) = \theta_i + \omega_i t + \frac{1}{2}\alpha t^2$
    \item $\omega_f^2 = \omega_i^2 + 2\alpha \Delta\theta$
    \item $\vec{a}_{total} = \vec{a}_c + \vec{a}_t$
\end{itemize}

\subsubsection*{Definiciones}
\begin{itemize}
    \item $\theta$: Posición angular [rad]
    \item $\omega$: Velocidad angular [rad/s]
    \item $\alpha$: Aceleración angular [rad/s$^2$]
    \item $r$: Radio [m]
    \item $T$: Periodo [s]
    \item $f$: Frecuencia [Hz]
\end{itemize}
\end{tcolorbox}

%===========================================================
% DINÁMICA LINEAL
%===========================================================
\begin{tcolorbox}[title=Dinámica Lineal]
\subsection*{Leyes de Newton}
\begin{itemize}
    \item \textbf{1ª Ley (Inercia):} Si $\sum \vec{F} = 0 \iff \vec{v} = \text{cte.}$
    \item \textbf{2ª Ley:} $\sum \vec{F} = m \cdot \vec{a}$
    \item \textbf{3ª Ley:} $\vec{F}_{A \to B} = - \vec{F}_{B \to A}$
\end{itemize}

\subsection*{Momentum e Impulso}
\begin{itemize}
    \item $\vec{p} = m \cdot \vec{v}$ (Momentum)
    \item $\vec{I} = \vec{F} \cdot \Delta t$ (Impulso, F cte.)
    \item $\vec{I}_{neto} = \Delta \vec{p} = \vec{p}_f - \vec{p}_i$
\end{itemize}

\subsection*{Fuerza Centrípeta}
\begin{itemize}
    \item $F_c = m \cdot a_c = m \frac{v^2}{r} = m \omega^2 r$
\end{itemize}

\subsubsection*{Definiciones}
\begin{itemize}
    \item $\vec{F}$: Fuerza [N]
    \item $m$: Masa [kg]
    \item $\vec{p}$: Momentum [kg·m/s]
    \item $\vec{I}$: Impulso [N·s]
\end{itemize}
\end{tcolorbox}


%===========================================================
% DINÁMICA ROTACIONAL
%===========================================================
\begin{tcolorbox}[title=Dinámica Rotacional]
\subsection*{Torque y Momento de Inercia}
\begin{itemize}
    \item $\vec{\tau} = \vec{r} \times \vec{F}$ (Torque)
    \item $|\tau| = r F \sin(\theta)$
    \item $I = \sum m_i r_i^2$ (Momento de Inercia)
\end{itemize}

\subsection*{Leyes de Rotación}
\begin{itemize}
    \item \textbf{2ª Ley (Rotación):} $\sum \vec{\tau} = I \cdot \vec{\alpha}$
\end{itemize}

\subsection*{Momentum Angular}
\begin{itemize}
    \item $\vec{L} = \vec{r} \times \vec{p}$ (partícula)
    \item $\vec{L} = I \cdot \vec{\omega}$ (cuerpo rígido)
    \item $\sum \vec{\tau} = \frac{d\vec{L}}{dt}$
\end{itemize}

\subsection*{Conservación del Momentum Angular}
\begin{itemize}
    \item Si $\sum \vec{\tau}_{ext} = 0 \implies \vec{L} = \text{cte.}$
    \item $I_i \omega_i = I_f \omega_f$
\end{itemize}

\subsubsection*{Definiciones}
\begin{itemize}
    \item $\vec{\tau}$: Torque [N·m]
    \item $I$: Momento de Inercia [kg·m$^2$]
    \item $\vec{L}$: Momentum Angular [kg·m$^2$/s]
\end{itemize}
\end{tcolorbox}

\end{multicols} % Finaliza el entorno de dos columnas

\end{document}
